% ============================================================
%  §7  Multi-scale / Multi-agent
% ============================================================

\section{Scale Inversion: Individual to Social}
\label{sec:multiscale}

\subsection{The Same Dynamics, Different Dominant Costs}

The dynamics described in preceding sections operate at the level of
an individual cognitive system.
At this scale, the dominant costs are recursion and compression:
the system must re-enter prior states at low cost, and must compress
high-dimensional continuous field states into reusable discrete handles.
Under these constraints, language tokens function as internal control
units.

When the same dynamics are extended to a population of interacting
agents, the dominant costs shift.
At the social scale, synchronization and coordination costs dominate:
tokens must stabilize shared trajectories across individuals by reducing
uncertainty about each other's landscape geometry.

The shift in dominant costs produces an apparent inversion of linguistic
function:
\begin{center}
\begin{tabular}{lll}
\toprule
\textbf{Scale} & \textbf{Dominant cost} & \textbf{Token function} \\
\midrule
Individual & Recursion, compression & Internal control units \\
Social     & Synchronization, coordination & Shared landscape anchors \\
\bottomrule
\end{tabular}
\end{center}

No property of the tokens has changed; only the effective constraint
structure changes.

\subsection{Social Conceptual Landscape}

As tokens circulate across a population, each agent's landscape is
perturbed by the same linguistic inputs.
When agents share sufficiently similar prior landscapes, repeated
shared perturbation produces correlated deepening of the same wells:
a shared conceptual landscape emerges.

The shared landscape is not stored in any individual agent; it is
distributed across the population and maintained through ongoing
linguistic interaction.
Its geometry is the product of all prior collective token--field
interactions.

\subsection{Norms, Institutions, and Paradigms}

Deeply sedimented regions of the social conceptual landscape correspond
to norms, institutions, and scientific paradigms.
These are not collections of meanings or propositions; they are stable
attractor regions of the collective $\Phi$ whose depth makes deviation
dynamically costly.

Institutional stability does not require enforcement by a central
authority; it follows from the depth of the collective potential well.
Institutional change requires coordinated perturbation sufficient to
overcome collective attractor depth --- a condition that accounts for
the historical difficulty of paradigm shifts.

\subsection{Intelligence as Landscape Maintenance}

At the social scale, intelligence is not a property of individual
agents but a distributed process:
\begin{quote}
\itshape
Intelligence is the socially distributed maintenance, synchronization,
and gradual reshaping of conceptual landscapes through language.
\end{quote}

Individual agents function as transient operators within this process:
they traverse the landscape, introduce local deformations, and transmit
stabilized structures to subsequent agents through externalized
linguistic logs.

The strict positivity $I>0$ means that the social landscape, like the
individual landscape, does not reach complete closure as an ordinary
dynamical state.
Further differentiation and innovation remain possible whenever new
perturbations or wider observation windows reactivate residual
asymmetries.
